% ============ 参考文献 ============
\chapter*{参考文献}
\addcontentsline{toc}{chapter}{参考文献}
\markboth{参考文献}{参考文献}

\begingroup
\zihao{5}
\setlength{\parindent}{0pt}
\linespread{1.3}\selectfont
\newcommand{\refitem}[2]{\noindent\makebox[8mm][l]{[#1]}\parbox[t]{\dimexpr\linewidth-8mm}{#2}\par\vspace{2pt}}

\refitem{1}{Meta Platforms. Ray-Ban Meta smart glasses[EB/OL]. (2023-09-27)[2026-04-10]. https://www.meta.com/smart-glasses/.}
\refitem{2}{Apple Inc. Apple Vision Pro[EB/OL]. (2024-02-02)[2026-04-10]. https://www.apple.com/apple-vision-pro/.}
\refitem{3}{KORAYEM M, TEMPLEMAN R, CHEN D, et al. ScreenAvoider: Protecting computer screens from ubiquitous cameras[EB/OL]. arXiv:1412.0008, 2014. https://arxiv.org/abs/1412.0008.}
\refitem{4}{TANG B, SHIN K G. Eye-Shield: Real-time protection of mobile device screen information from shoulder surfing[C]//32nd USENIX Security Symposium. 2023. arXiv:2308.03868. https://arxiv.org/abs/2308.03868.}
\refitem{5}{GE S, XIA Z, TONG Y, et al. A screen-shooting resilient document image watermarking scheme using deep neural network[EB/OL]. arXiv:2203.05198, 2022. https://arxiv.org/abs/2203.05198.}
\refitem{6}{GOODFELLOW I J, SHLENS J, SZEGEDY C. Explaining and harnessing adversarial examples[C]//International Conference on Learning Representations (ICLR). 2015.}
\refitem{7}{MADRY A, MAKELOV A, SCHMIDT L, et al. Towards deep learning models resistant to adversarial attacks[C]//International Conference on Learning Representations (ICLR). 2018.}
\refitem{8}{CARLINI N, WAGNER D. Towards evaluating the robustness of neural networks[C]//2017 IEEE Symposium on Security and Privacy (SP). Piscataway: IEEE, 2017: 39-57.}
\refitem{9}{AKHTAR N, MIAN A. Threat of adversarial attacks on deep learning in computer vision: a survey[J]. IEEE Access, 2018, 6: 14410-14430.}
\refitem{10}{BERNSTEIN D J. ChaCha, a variant of Salsa20[R]. Chicago: The University of Illinois at Chicago, 2008.}
\refitem{11}{KRAWCZYK H, ERONEN P. HMAC-based extract-and-expand key derivation function (HKDF)[S]. RFC 5869. Internet Engineering Task Force, 2010.}
\refitem{12}{JOCHER G, CHAURASIA A, QIU J. Ultralytics YOLOv8[EB/OL]. (2023-01-10)[2026-04-10]. https://github.com/ultralytics/ultralytics.}
\refitem{13}{SMITH R. An overview of the Tesseract OCR engine[C]//Ninth International Conference on Document Analysis and Recognition (ICDAR). Piscataway: IEEE, 2007: 629-633.}
\refitem{14}{SHARMA G, WU W, DALAL E N. The CIEDE2000 color-difference formula: implementation notes, supplementary test data, and mathematical observations[J]. Color Research \& Application, 2005, 30(1): 21-30.}
\refitem{15}{IEEE. IEEE recommended practices for modulating current in high-brightness LEDs for mitigating health risks to viewers: IEEE Std 1789-2015[S]. New York: IEEE, 2015.}
\refitem{16}{SZEGEDY C, ZAREMBA W, SUTSKEVER I, et al. Intriguing properties of neural networks[C]//International Conference on Learning Representations (ICLR). 2014.}
\refitem{17}{DODIS Y, RISTENPART T, STEINBERGER J, et al. To hash or not to hash again? (in)differentiability results for H\textsuperscript{2} and HMAC[C]//Advances in Cryptology (CRYPTO). Berlin: Springer, 2012: 348-366.}
\refitem{18}{TRAN V, JAYATILAKA G, ASHOK A, et al. DeepLight: Robust and unobtrusive real-time screen-camera communication for real-world displays[C]//ACM/IEEE International Conference on Information Processing in Sensor Networks (IPSN). 2021. DOI: 10.1145/3412382.3458269.}
\refitem{19}{NISHAR A A M, KUDEKAR S, KINTZING B, et al. Revelio: A real-world screen-camera communication system with visually imperceptible data embedding[C]//IEEE International Conference on Acoustics, Speech and Signal Processing (ICASSP). 2025. arXiv:2501.02349. https://arxiv.org/abs/2501.02349.}
\refitem{20}{GURI M, BYKHOVSKY D, ELOVICI Y. BRIGHTNESS: Leaking sensitive data from air-gapped workstations via screen brightness[C]//2019 12th CMI Conference on Cybersecurity and Privacy. Piscataway: IEEE, 2019. DOI: 10.1109/CMI48017.2019.8962137.}

\endgroup
